\chapter*{Glosario}

\begin{longtable}{>{\bfseries\raggedright\arraybackslash}p{3.4cm}L{10.6cm}}
Actas & Documentos de seguimiento de reuniones utilizados para reconstruir decisiones de alcance, arquitectura y planificación.\\
AI Act & Reglamento Europeo de Inteligencia Artificial, considerado en la memoria como marco normativo relevante para posibles fases futuras con uso sanitario o de rehabilitación.\\
Algoritmo evolutivo & Método de optimización basado en poblaciones, selección y variación de soluciones candidatas. En el proyecto se usa para ajustar los pesos de la red neuronal.\\
Aprendizaje evolutivo & Proceso por el que una población de vehículos conserva, muta y reevalúa modelos de control a partir de su fitness.\\
ASPAYM & Entidad colaboradora vinculada al contexto de rehabilitación y valoración funcional para el que se plantea el simulador.\\
Blueprint & Sistema visual de programación de Unreal Engine empleado para implementar lógica de juego, interacción, controladores y sistemas de simulación.\\
Ceda el paso & Señal de prioridad que permite continuar si el cruce está libre, pero exige esperar cuando existen vehículos con preferencia.\\
Chaos Vehicle & Sistema de físicas de Unreal Engine 5 utilizado para simular el comportamiento dinámico del vehículo.\\
CSV & Formato de texto tabular usado para exportar métricas de entrenamiento, test y depuración.\\
DCA & Daño Cerebral Adquirido.\\
Fitness & Función de aptitud que cuantifica el comportamiento de cada vehículo autónomo durante una generación evolutiva.\\
FreeRun & Modo de simulación libre con respawn y tráfico continuo, usado como diagnóstico de convivencia entre vehículos y no como test generacional comparable.\\
Git & Sistema de control de versiones utilizado para registrar cambios del proyecto.\\
GitLab & Plataforma colaborativa empleada para issues, milestones y trazabilidad de tareas cerradas.\\
IA & Inteligencia artificial. En esta memoria se refiere principalmente al controlador neuronal y al entrenamiento evolutivo de vehículos.\\
Inferencia & Cálculo de las salidas de la red neuronal a partir de las entradas sensoriales y contextuales del vehículo.\\
Intersección & Zona del mapa donde el vehículo debe decidir trayectoria, prioridad y respuesta ante señalización o tráfico.\\
IT & Identificador abreviado de triggers o actores de intersección utilizado en los informes de diagnóstico, por ejemplo \texttt{IT58}.\\
LOPDGDD & Ley Orgánica 3/2018 de Protección de Datos Personales y garantía de los derechos digitales.\\
NavMesh & Malla de navegación utilizada por motores de videojuegos para calcular rutas transitables. En este proyecto fue estudiada y posteriormente descartada como solución principal para vehículos físicos.\\
Neuroevolución & Entrenamiento de redes neuronales mediante algoritmos evolutivos. En este trabajo se evoluciona el conjunto de pesos, manteniendo fija la topología.\\
NPC & \textit{Non-Player Character}; agente controlado por el sistema y no por el usuario.\\
Path Finding & Búsqueda de caminos entre un origen y un destino. Se evaluó como alternativa inicial, pero no resolvía por sí sola el control físico local del coche.\\
Perceptrón multicapa & Red neuronal con una capa de entrada, una o varias capas ocultas y una capa de salida.\\
Raycasting & Técnica de consulta geométrica que lanza rayos o trazas desde un actor para detectar obstáculos y distancias en el entorno.\\
Red neuronal & Modelo computacional formado por pesos y funciones de activación que transforma entradas en salidas de control.\\
RF & Requisito funcional; describe una capacidad que el sistema debe proporcionar.\\
RGPD & Reglamento General de Protección de Datos de la Unión Europea.\\
RNF & Requisito no funcional; describe una propiedad de calidad, trazabilidad, rendimiento o seguridad del sistema.\\
Rotonda & Intersección circular con entradas y salidas que exige gestionar prioridad, velocidad, trayectoria y salida prevista.\\
Sensor & Traza simulada o variable contextual que alimenta el vector de entrada del controlador neuronal.\\
Simulación urbana & Entorno virtual que reproduce una ciudad con calzada, señales, intersecciones, vehículos y eventos de tráfico.\\
STOP & Señal que exige parada explícita antes de continuar. En el proyecto se valida mediante posición legal, espera mínima y estado del cruce.\\
Test & Ejecución con modelo fijo y mutación desactivada, usada para evaluar comportamiento en puntos y trayectorias diferenciados.\\
Train/test & Separación entre sesiones de entrenamiento, que modifican modelos, y sesiones de test, que evalúan pesos ya fijados.\\
Unreal Engine 5 & Motor de desarrollo 3D utilizado como base del simulador.\\
Validación experimental & Comprobación mediante ejecuciones, CSV e informes de análisis. No equivale a validación clínica con pacientes.\\
Vehículo autónomo & Coche controlado por el sistema mediante sensores, red neuronal, señales de tráfico y gestor evolutivo.\\
WCAG & \textit{Web Content Accessibility Guidelines}; guía de accesibilidad citada como referencia parcial para futuras interfaces digitales del simulador.\\
\end{longtable}
